\chapter{The filled hull}\label{ch:hull}
\module{NoCompromise.Hull.Defs,
        NoCompromise.Hull.Properties,
        NoCompromise.Stationary.Defs}

\begin{notation}\label{not:stationary}
\uses{def:stationary}
Throughout Part~\ref{part:capacity}, $\Omega$ is a bounded connected
$C^{3,\beta}$ stationary domain of volume $V$ with multiplier $\lambda$
(Definition~\ref{def:stationary}), as produced by
Corollary~\ref{cor:minimizer-stationary}.
\end{notation}

\begin{definition}[Filled hull]\label{def:hull}
\uses{not:stationary}
Let $U_\infty$ be the unbounded connected component of
$\R^3\setminus\overline\Omega$ and put
\begin{equation}\label{eq:hull}
  K:=\R^3\setminus U_\infty .
\end{equation}
Thus $K$ fills every bounded complementary component of $\Omega$.
\end{definition}

\begin{lemma}[Hull properties]\label{lem:hull-properties}
\uses{def:hull,thm:component-count}
$K$ is compact and regular closed, $K=\overline{\operatorname{int}K}$; the sets
$\operatorname{int}K$ and $\R^3\setminus K$ are connected; $\partial K$ is
connected; and
\begin{equation}\label{eq:hull-inclusions}
  \overline\Omega\subset K,\qquad\partial K\subset\partial\Omega .
\end{equation}
Moreover $\partial K$ is $C^{3,\beta}$, the outward normals and curvatures of
$K$ and $\Omega$ agree on $\partial K$, and
\begin{equation}\label{eq:hull-perimeter}
  \Per(K)\le\Per(\Omega).
\end{equation}
\end{lemma}

\begin{proof}\uses{thm:component-count,prop:no-singular-points,thm:tubular,
                    prop:connected}
Write
$\partial\Omega=\Sigma_1\mathbin{\dot\cup}\cdots\mathbin{\dot\cup}\Sigma_k$
(Proposition~\ref{prop:no-singular-points}).
Theorem~\ref{thm:component-count} says $\R^3\setminus\partial\Omega$ has
$k+1$ components.  One of them is the connected set $\Omega$
(Proposition~\ref{prop:connected}), so there are exactly $k$ components of
$\R^3\setminus\overline\Omega$.  Every $\Sigma_i$ has one collar side in
$\Omega$ and one in a complementary component
(Theorem~\ref{thm:tubular}); every complementary component has nonempty
boundary contained in $\partial\Omega$, so it meets at least one $\Sigma_i$.
Hence the map sending $\Sigma_i$ to the component on its non-$\Omega$ side is
a surjection between finite sets of size $k$, therefore a bijection.  Exactly
one boundary component faces $U_\infty$; call it $\Sigma_\infty$.

By construction $\partial K=\partial U_\infty=\Sigma_\infty$; all other
boundary components become interior interfaces once their adjacent cavities are
filled.  The union of $\Omega$, all bounded complementary components, and
collars crossing their common interfaces is connected and equals
$\operatorname{int}K$, and the outer collar gives
$K=\overline{\operatorname{int}K}$.  Normal agreement holds because both
outward normals point into $U_\infty$.  Finally
$\Per(K)=\Hs^2(\Sigma_\infty)\le\sum_i\Hs^2(\Sigma_i)=\Per(\Omega)$.
\end{proof}

\begin{fnote}[Representative bookkeeping]\label{fnote:bookkeeping}
Five distinct objects are now in play: the measurable set $E$, its density-one
open representative $\Omega$, its closure $\overline\Omega$, the filled hull
$K$, and the exterior $\R^3\setminus K$.  Confusing any two of them produces a
false statement.  A Lean development should give each a distinct name and
prove the inclusions \eqref{eq:hull-inclusions} once, as a reusable lemma,
rather than re-deriving them in each consumer.
\end{fnote}

